%% notes-tex/019-simb-multimodal/sections/2-expression.tex
%% [[experiments.019-simb-multimodal.expression-round-retrospective]]
%% SOURCE: numbers are from experiments/019-simb-multimodal/results/*.json (committed
%% artifacts) and from named W&B run ids in project zhao-group/torchcell_019_expr_v8.
%% Nothing here is hand-authored.

\section{Expression\secstatus{todo}}
\label{sec:expression}

\subsection{The task and its ceiling\secstatus{todo}}

The \term{expression strand} predicts, for a strain carrying one gene deletion, the
$\log_2$ ratio of every measured gene's transcript abundance against wild type. Targets
come from the Kemmeren 2014 knockout microarray compendium (doi:10.1016/j.cell.2014.02.054)
and the Sameith 2015 double-deletion compendium (doi:10.1038/sdata.2015.15), fused in the
\file{fig3_core} build. Of the 1,556 records in that build, 1,484 are single deletions and
72 are Sameith doubles, so $95.4\,\%$ of the training signal carries exactly one perturbed
gene.
\src{experiments/019-simb-multimodal/results/fig3_overlap_census.json}

The \term{reproducibility ceiling} is the highest correlation any predictor could reach
against a target that carries measurement noise. It is estimated here from the 82 deletions
measured independently in both compendia, as the mean over reporter genes of
$\sqrt{\operatorname{clip}(\hat\rho_g, 0, 1)}$ where $\hat\rho_g$ is that gene's test-retest
correlation across the paired strains. The estimate is $\mathbf{0.7746}$.
\src{experiments/019-simb-multimodal/results/expression_ceiling_replicate.json}

The best validation score observed on this task is
\file{pearson_per_feature} $= \mathbf{0.2407}$, from run \texttt{hx8pxdic} in project
\file{torchcell_019_expr_v9}, whose peak sits at epoch 9,188 of the 9,997 it reached.
The best in \file{torchcell_019_expr_v8} is $0.2293$, run \texttt{b50f93ju}, peaking at
epoch 4,007 of 4,105. Both are $\approx 30\,\%$ of the ceiling, and both peak in the last
tenth of their run.
\src{experiments/019-simb-multimodal/results/round_leaderboards.csv}

The metric is the per-reporter correlation across strains, averaged over reporters, which
is the quantity that goes to zero when a model predicts each gene's mean; the companion
per-strain metric stays near $0.4$ under exactly that collapse and is not the objective.
The scoring rule is \file{roll_max}, the maximum of a centered five-point rolling mean of
the validation curve. It is an upward-biased order statistic whose bias grows with the
number of epochs run, so it is comparable only between runs of similar length, which is why
the epoch count travels with every number quoted from the leaderboard.

The manuscript's Results section currently states $r = 0.543$ for expression. That string
sits inside a paragraph marked \texttt{[FILLER -- R3.]} and is a placeholder, not a
measurement.

\subsection{The finding of the round: nothing was trained to saturation\secstatus{todo}}

Sixty-seven runs across waves 1 to 4b were scored, and not one of them ran past 300 epochs.
Against a validation curve that dips at epochs 200 to 300 and then climbs past its early
peak, every one of those runs was stopped inside the dip.
\src{experiments/019-simb-multimodal/results/decoder_arms_torchcell_019_expr_v8.csv}

\begin{table}[htbp]
\centering\small
%% SOURCE: recomputed in the round retrospective from decoder_arms_torchcell_019_expr_v8.csv
\begin{tabular}{lrl}
\toprule
wave & runs & epoch range \\
\midrule
wave 1  & 10 & 70--140 \\
wave 2  &  5 & 76--140 \\
wave 3  & 36 & 18--276 \\
wave 4a &  8 & 77--81 \\
wave 4b &  8 & 300 \\
\midrule
\textbf{all} & \textbf{67} & \textbf{18--300, none above 300} \\
\bottomrule
\end{tabular}
\caption[]{Every scored arm in the expression round stopped at or below 300 epochs.}
\label{tab:expr-truncation}
\end{table}

Four long runs were later allowed to continue. All four died at 12.60 to 12.61 hours of
wall clock at 1,620 to 1,625 epochs, which is a scheduler limit rather than convergence,
and their validation peaks sit at 93 to $99\,\%$ of the way through the run. Eval-mode
training Pearson reached $0.652$ to $0.725$ and was still climbing when they were killed.

Longer runs have since been allowed. The picture did not change: the best v8 run peaks at
epoch 4,007 of 4,105 and the best v9 run at epoch 9,188 of 9,997. Raising the epoch budget
from 1,600 to 10,000 moved the score from $0.204$ to $0.241$ and moved the peak with it, so
the curve is still rising where it was cut off. \textbf{No peak has ever been observed on
this task}, at any budget tried, and the saturation epoch remains unknown.

\notesec{Consequence} An arm comparison at 300 epochs is not a small effect measured
noisily. Both arms were stopped before either had the chance to separate, so the difference
between them reflects early-training transients. The null sink arm reads
$+0.0024 \pm 0.0090$ over four seed pairs and post-perturbation mixing reads
$-0.0023 \pm 0.0076$, and neither number is evidence about its mechanism. More seeds do not
repair this: replicates buy precision on the wrong estimand.

\subsection{Loss and metric move in opposite directions\secstatus{todo}}

Validation \texttt{expression/loss} reaches its minimum long before validation Pearson
peaks, then rises for the remaining $\approx 1{,}200$ epochs while Pearson keeps climbing.

\begin{table}[htbp]
\centering\small
%% SOURCE: experiments/019-simb-multimodal/results/wave_history_stage-wave4b.csv
\begin{tabular}{lll}
\toprule
run & val loss min & val Pearson max \\
\midrule
\file{N_sink_mm} s1  & 0.2640 @ ep 53  & 0.0640 @ ep 257 \\
\file{N_sink_mm} s42 & 0.2653 @ ep 217 & 0.0824 @ ep 296 \\
\file{R_ref} s1       & 0.2639 @ ep 42  & 0.0632 @ ep 293 \\
\file{R_ref} s42      & 0.2662 @ ep 202 & 0.0701 @ ep 297 \\
\bottomrule
\end{tabular}
\caption[]{Minimum validation loss precedes maximum validation Pearson by 40 to 250 epochs
in wave 4b, and by $\approx 1{,}200$ epochs in the long runs. The 010 fitness and
interaction task does not do this.}
\label{tab:expr-divergence}
\end{table}

Two things follow. Every checkpoint the expression strand ever saved by validation loss was
the early-dip model, roughly 1,300 epochs before the good one, which is why
best-by-metric checkpointing was added. And the divergence itself is information about the
objective: a quantile loss on a $\log_2$ ratio is minimized by a well-calibrated, narrow
predictive distribution, while Pearson rewards getting the ordering right at whatever
scale. The prediction standard-deviation ratio moving from $0.001$ to $0.52$ over the same
window is consistent with the model widening its predictions long after the loss says it
should stop.

\notesec{What this does not license} The mechanism sentence above is a reading of the
curves, not a measurement. Nothing has tested whether a different loss recovers the
Pearson gain earlier. That is one of the options in \cref{sec:directions}.

\subsection{The perturbation operator has no pair term\secstatus{todo}}

At $|S_b| = 1$ the cross-attention that carries the perturbation runs its softmax over a
one-element key set, so the attention weight is identically 1. Re-drawing $W_Q$ and $W_K$
at standard deviation 10 changes the output by exactly $0.0$, and 16,200 of 32,760
attention parameters are dead.
\src{experiments/019-simb-multimodal/results/perturbation_selector_degeneracy.json}

The model therefore reduces to $g(h_i + c_b)$: gene identity $i$ and strain identity $b$
meet once, by addition, and no term depends on the pair $(p, i)$ of deleted gene and
reporter gene. The rank of the pair term each candidate mechanism can express is what
wave 6 was built to ladder.

\begin{table}[htbp]
\centering\small
%% SOURCE: experiments.019-simb-multimodal.multiplicative-perturbation-conditioning (derivation)
\begin{tabular}{lll}
\toprule
mechanism & form & pair rank \\
\midrule
\file{V_ref} additive & $g(h_i + c_b)$ & \textbf{0} \\
\file{V_sink} gated attention & one bounded scalar per head & 9 \\
\file{V_basis}$\{16,32,64\}$ & $\sum_j b_{ij}(h_i)\,a_j(z_S)$ & $r$ \\
\file{V_film}, \file{V_hadamard} & diagonal multiplicative & $d = 90$ \\
\bottomrule
\end{tabular}
\caption[]{Pair-rank ladder. The reference arm cannot express any dependence on which
reporter is being predicted given which gene was deleted.}
\label{tab:pair-rank}
\end{table}

Capacity is not what binds. A rank-$r$ reconstruction of the target matrix reaches $0.7265$
at $r = 32$ and $0.7799$ at $r = 64$, against a best model score of $0.241$, so a null on
this axis is a statement about the mechanism.
\src{experiments/019-simb-multimodal/results/lowrank_output_ceiling.json}

\subsection{The largest signal in the strand is imputation, and it is orthogonal\secstatus{todo}}

A ridge oracle that observes $m$ randomly chosen genes of a held-out strain and predicts
the rest reaches far past anything the model does.

\begin{table}[htbp]
\centering\small
%% SOURCE: experiments/019-simb-multimodal/results/masked_conditioning_oracle.json
%% SOURCE: experiments/019-simb-multimodal/results/cross_study_conditioning_oracle.json
\begin{tabular}{lrrrr}
\toprule
$m$ observed & within-Kemmeren & within-Sameith & Kem $\to$ Sam & Sam $\to$ Kem \\
\midrule
10   & 0.4562 & 0.3649 & 0.2335 & 0.2122 \\
100  & 0.6693 & 0.6417 & 0.3815 & 0.3922 \\
1000 & 0.7832 & 0.7779 & \textbf{0.4838} & \textbf{0.4803} \\
\bottomrule
\end{tabular}
\caption[]{Masked-conditioning oracle, \file{pearson_per_feature} on validation
strains. The within-study value at $m = 1000$ exceeds the $0.775$ reproducibility ceiling,
which the cross-study column resolves: same-array conditioning was predicting the target's
own measurement noise.}
\label{tab:oracle}
\end{table}

The gain is not a better genotype-to-expression map. Removing a genotype predictor, a
$k$-nearest-neighbor model on \file{prot_T5_all} with $k = 25$, leaves $97.5$ to
$100.6\,\%$ of the conditioning gain intact, so the channel the objective optimizes is
switched off at $m = 0$, which is where the strand is scored.
\src{experiments/019-simb-multimodal/results/conditioning_gain_after_genotype.json}

\tcfigfit{figures/residual_covariance_diagnostic.pdf}{Residual gene-gene structure in the
expression compendium after removing each gene's mean. The pattern replicates on held-out
strains and lives in about 33 effective components, which is what a per-gene independent
readout discards.}{fig:residual-cov}

What makes the oracle worth keeping is the structure it exploits. After removing each
gene's mean, the residual gene-gene correlation pattern replicates on held-out strains at
split-half $r = 0.8687$ against a permutation null of $8.45\times10^{-5}$, and it has an
effective rank of $32.78$, with $59.1\,\%$ of variance inside the top 32 components. A
per-gene independent readout emits 6,127 marginals and discards all of it.
\src{experiments/019-simb-multimodal/results/residual_covariance_diagnostic.json}

\subsection{Why the byte count is misleading\secstatus{todo}}

The expression compendium holds more scalar labels than the trigenic interaction data the
model does well on, which is the intuition that predicted this task would be easy. The
label count is not the axis that matters.

\begin{table}[htbp]
\centering\small
%% SOURCE: row counts read from $DATA_ROOT/data/torchcell/<dataset>/preprocess/data.csv
\begin{tabular}{lrrrp{42mm}}
\toprule
dataset & records & per record & scalars & perturbation axis \\
\midrule
Kuzmin 2018 trigenic interaction  & 91{,}111 & 1       & 91{,}111    & triples; a gene recurs with many partners \\
Kemmeren 2014 knockout expression & 1{,}484  & 6{,}169 & 9{,}153{,}196 & 1{,}484 single deletions, each once \\
Ohya 2005 morphology              & 4{,}718  & 281     & 1{,}325{,}758 & 4{,}718 single deletions, each once \\
\bottomrule
\end{tabular}
\caption[]{Scalar labels are not the binding resource. Expression carries $\approx 9.2$
million scalars against the trigenic set's 91 thousand, but its perturbation axis has 1,484
points and no gene is deleted twice, so nothing in the data isolates the effect of a
deletion from the effect of the one array it was measured on.}
\label{tab:label-accounting}
\end{table}

In the trigenic data a given gene appears in many records paired with different partners,
so the supervision repeatedly constrains that gene's contribution. In the expression data a
gene appears as the perturbation exactly once. The 6,169 labels attached to it are 6,169
readouts of a single perturbation event, correlated with each other at the level the
residual covariance above measures, so they are worth far less than 6,169 independent
constraints on the deletion.

\notesec{Established} The expression strand sits at $0.241$ against a ceiling of $0.775$.
The round produced no valid comparison between mechanisms, because every arm was scored
before saturation. It produced three things that hold: the additive operator provably
cannot express a pair term; loss and metric diverge on this target and not on 010; and the
one large available signal is imputation from other genes of the same strain, measured to
be orthogonal to the genotype-only task.

\notesec{Cannot say} Whether any of the pair-rank mechanisms help. Whether the graph prior
is the right prior for this task. Whether a different objective closes the loss-metric gap.
